% Section 4.1, from lectures/L04/appendix/a_fifo.md.

\section{Designing \code{fifo.vhd}}
\label{c4:sec:fifo}\label{c4:app:a}
The datapath produces and consumes bytes at different rates:
\begin{itemize}
  \item Software writes a byte to send whenever it likes.
  \item The transmitter sends one every ten bit-times.
  \item The receiver delivers a byte every frame.
  \item Software reads it whenever it gets round to it.
\end{itemize}

A \textbf{FIFO} (a first-in, first-out queue) absorbs that mismatch: bytes go in one end and come
out the other in the same order, and each side can run without waiting for the other, up to the
depth of the queue.

The register bank (\secref{c5:app:a}) needs two of these, one for each direction, so \code{fifo} is
written once as a generic module and instantiated twice.

\subsection{Interface}
\label{c4:sec:fifo_interface}

\bookfigure{fifo}{Module \code{fifo}}{c4:fig:fifo}

\begin{booktable}{@{}p{6.4em}p{3.2em}p{17.6em}L@{}}
  \thead{Port / generic} & \thead{Dir} & \thead{Type} & \thead{Meaning} \\
  \midrule
  \code{WIDTH} & generic & \code{natural range 1 to 64 := 8} & Bits per entry. \\
  \code{DEPTH} & generic & \code{natural range 1 to 256 := 8} & Number of entries. \\
  \code{clock} & in & \code{std_logic} & 50\,MHz system clock. \\
  \code{reset_s2_n} & in & \code{std_logic} & Active-low synchronized reset. \\
  \code{wdata} & in & \code{std_logic_vector(WIDTH-1 downto 0)} & The entry to push. \\
  \code{wr} & in & \code{std_logic} & \code{'1'} to push \code{wdata} this cycle. \\
  \code{rd} & in & \code{std_logic} & \code{'1'} to pop the front entry this cycle. \\
  \code{rdata} & out & \code{std_logic_vector(WIDTH-1 downto 0)} & The entry at the front
    (combinational). \\
  \code{empty} & out & \code{std_logic} & \code{'1'} if the FIFO is empty. \\
  \code{full} & out & \code{std_logic} & \code{'1'} if the FIFO is full. \\
\end{booktable}

\code{fifo_tb} binds the ports positionally and drives the generic map too, instantiating a depth-4
FIFO. \code{rdata} shows the front entry continuously, without a \code{rd}; \code{rd} is what
\emph{discards} it and moves to the next. That split, look then advance, is exactly what the RX path
in \chapref{5} needs.

\subsection{Behaviour}
\label{c4:sec:fifo_behaviour}
A ring buffer: an array of \code{DEPTH} entries, a write pointer (\code{head}), a read pointer
(\code{tail}), and a count. Each pointer wraps around from \code{DEPTH-1} back to 0.

On reset (\code{reset_s2_n} low), \code{head}, \code{tail} and \code{count} all go to 0, so the FIFO
comes out of reset empty and not full; as everywhere else in this peripheral the reset is asserted
asynchronously and tested ahead of \code{rising_edge(clock)}. The entries themselves need no
clearing, since with \code{count = 0} nothing in the array is reachable and the first write lands on
entry 0 anyway. What does matter is that the two pointers start \emph{equal}, not merely at some
known value: \code{empty} and \code{full} are read off \code{count}, so a \code{head} and
\code{tail} that disagreed out of reset would keep the count honest while handing back a stale entry
instead of the byte just pushed. \code{rdata} still shows an entry while the FIFO is empty (entry 0
out of reset, a stale one after a drain), whatever that entry happens to hold, which is one more
reason the caller watches \code{empty} rather than trusting \code{rdata}.

On each rising edge, while not in reset, a \textbf{write} happens only if \code{wr = '1'} and the
FIFO is not full, storing \code{wdata} at \code{head} and advancing \code{head}; a \textbf{read}
happens only if \code{rd = '1'} and the FIFO is not empty, advancing \code{tail}. The count follows
from the pair: it rises by one on a write without a read, falls by one on a read without a write,
and is unchanged when both or neither happen.

\code{full} is \code{count = DEPTH}, \code{empty} is \code{count = 0}, and \code{rdata} is
combinationally the entry at \code{tail}. Guarding writes on \code{not full} and reads on
\code{not empty} means the pointers can never cross; a write to a full FIFO and a read from an empty
one are simply dropped, which is why the caller must watch the flags.

\subsection{What the testbench pins down}
\label{c4:sec:fifo_tb}
\code{fifo_tb} uses a depth-4 instance. It checks \code{empty} after reset and \code{full} after
four pushes, then pushes a fifth entry while full and confirms it is dropped rather than stored over
an existing one, and finally drains the FIFO and checks the four bytes come back in the order they
went in, leaving it \code{empty} again. A last case raises \code{wr} and \code{rd} on the same edge
with two entries queued: one entry in and one out, so the depth must not move, and the byte pushed
alongside the pop must still be there after the next one. That is the case where a count adjusted
in two independent branches quietly loses an entry.

That is the whole contract: order preserved, and the flags honest at both ends.

\subsection{Where it fits}
\label{c4:sec:fifo_fits}
The register bank (\secref{c5:app:a}) instantiates two \code{fifo}s, both \code{WIDTH = 8, DEPTH =
8}: a TX FIFO that a \code{TX_DATA} register write pushes and the transmitter drains, and an RX FIFO
that the receiver pushes and an \code{RX_DATA} read plus an \code{RX_POP} write drains. The FIFO's
\code{full} and \code{empty} flags become the \code{STATUS} register's TX-ready and RX-valid bits.

\subsection{What's ahead}
The exercises in \secref{c4:sec:exercises} build \code{fifo} and run its testbench, put the receive
path of \chapref{3} into \code{uart_top}, and work out why overrun is not the receiver's question to
answer. Then \secref{c5:app:a} builds \code{uart_regs}, the register bank, which wraps two of these
FIFOs and the datapath in the register map that software drives.
